\scanpage{22}\mspage{21}
\begingroup\csname tagsleft@true\endcsname

\unsure{termes de|lemme de} $x\in\operatorname{omb}(X')\subset\operatorname{omb}(X\unsure{.})$, $y\in\operatorname{om}(Y')\subset\operatorname{omb}(Y)$ par $x\mathrel{\Vert}y$. Or la condition $X'_L=\varnothing$ implique $x_L=\varnothing$ car $x_L\ll X'_L$, et \deleted{\unsure{ce|elle}}\unsure{cela} signifie \deleted{que} ici que $\unsure{x|x'}\notin\operatorname{omb}(L)$ (cf lemme) et idem pour $y$.

Donc $x\in\operatorname{omb}(X)\setminus\operatorname{omb}(L)$, $y\in\operatorname{omb}(Y)\setminus\operatorname{omb}(L)$ et si on a At CL(spac), on conclut bien $x\mathrel{\Vert}y$, \deleted{\unsure{d'où}} Donc la condition
\[
 \forall x\in\operatorname{omb}(X)\setminus\operatorname{omb}(\deleted{\unsure{X|Y}}L),\
 y\in\operatorname{omb}(Y)\setminus\operatorname{omb}(L),\
 \text{\deleted{on} a}\quad x\mathrel{\Vert}y
 \tag*{(AL)}
\]
implique
\[
 \forall \unsure{X|X'}\in\operatorname{Omb}(\unsure{x|X|x'}),\
 Y'\in\operatorname{Omb}(Y),\ X'_L=Y'_L=\varnothing,\ \text{on a}\quad X\mathrel{\Vert}Y
 \tag*{(A)}
\]
et l'inverse est trivial, comme on voit en prenant $X,Y\in\mathcal{L}$. On a utilisé ici le

\uline{Lemme 1} \unsure{Soit|Soient} $\deleted{\unsure{L|E}\leq F}\ P\leq Q$ (en l'occurrence $\deleted{\unsure{X}}L\leq X$ ou $L\leq Y$) $x\in\unsure{\operatorname{omb}|\operatorname{Omb}}(Q)$, alors $\deleted{x\mathrel{\Vert}Q\unsure{\Rightarrow|-}}$
\[
 x_P=\varnothing\Longleftrightarrow x\notin\operatorname{omb}(\deleted{\unsure{Q}}P)\quad
 \text{i.e}\quad x\not\ll\deleted{\unsure{Q}}P\quad
 (\deleted{\unsure{\Longleftrightarrow}x\mathrel{\Vert}\unsure{P|Q}})
\]

En effet, $\deleted{x\ll P\Longleftrightarrow x_P\unsure{\Rightarrow|=x}}$ pour la première équivalence, passant aux négations on \deleted{\unsure{trouve}}\inserted{doit} prouver $x\ll P\Longleftrightarrow x_P\ne\varnothing_{\unsure{\mathscr{F}|F|7}}$, et comme $x_P\ll x$ et $x$ minimal dans $M\setminus\{\varnothing_{\mathscr{F}}\}$ pour $\ll$, $x_P\ne\varnothing_{\mathscr{F}}$ signifie aussi $x_P=x$ i.e. $x\ll P$ OK.

\deleted{D'autre part, $x\mathrel{\Vert}P\Rightarrow x\not\ll P$ est trivial, \unsure{vu} la définition} Je profite de l'occasion pour expliciter

\noindent\vrule\hspace{.5em}\begin{minipage}{.94\linewidth}
\uline{Lemme 2} \inserted{(moyennant At D)} Soient $P\leq Q$, et $F\ll Q$. Conditions équivalentes
\[
 F_P=\varnothing\Longleftrightarrow F\mathrel{\Vert}P
\]
\end{minipage}

On applique At D (ou plutôt, la prop. qui suit, \unsure{haut|haut de} page \unsure{15|16}) à $P,Q$ $P\ll P$, $F\ll Q$, pour avoir $|\Rightarrow$ ; l'autre est vrai avec les seuls axiomes At 1 -- At 7
\endgroup
